\section{半正定规划优化矩阵不等式}
\label{sec:09-Convex-Optimization/04-Convex-Problems/04}

两百多个服务要分到两个机房，跨机房的每一次调用都要付一份网络延迟。每个服务归哪边，是一个二值决策 \(x_i\in\set{-1,+1}\)；总代价是所有跨边调用量的和。你写了个贪心，跑出一个方案，成本 \(1.82\) 亿次调用。运维问：这个方案离最优还有多远。

穷举需要看 \(2^{240}\) 种划分，这个数比可观测宇宙的原子还多。再跑十次随机重启，最好的一次是 \(1.80\)，也仍然只是``我没找到更好的''。

有一条路能给出别的东西：把这个离散问题松弛成一个半正定规划，解出来是 \(1.71\)。这个数不是某个方案的成本，而是一句关于全部 \(2^{240}\) 种划分的断言——没有任何一种能低于它。于是那份贪心方案的评语可以改写成：至多比最优差 \(6\%\)。差多少是算出来的，不是猜的。

\begin{center}
\begin{tikzpicture}[x=1cm,y=1cm,font=\small,>=stealth]
  \fill[black!6] (2.3,2.0) ellipse (2.6 and 2.3);
  \draw[black!50,thick] (2.3,2.0) ellipse (2.6 and 2.3);
  \fill[black!14] (1.1,1.0) -- (2.2,0.4) -- (3.6,1.1) -- (4.0,2.6) -- (3.0,3.5) -- (1.5,3.2) -- (0.5,2.0) -- cycle;
  \draw[black!70,dashed] (1.1,1.0) -- (2.2,0.4) -- (3.6,1.1) -- (4.0,2.6) -- (3.0,3.5) -- (1.5,3.2) -- (0.5,2.0) -- cycle;
  \foreach \a/\b in {1.1/1.0,2.2/0.4,3.6/1.1,4.0/2.6,3.0/3.5,1.5/3.2,0.5/2.0,2.3/1.9}
    \fill[black] (\a,\b) circle (2pt);
  \draw[dashed,black!60] (-0.3,1.129) -- (1.129,-0.3);
  \draw[dashed,black!60] (-0.3,2.4) -- (2.4,-0.3);
  \draw[black,thick] (0.353,0.476) circle (2.8pt);
  \draw[black,thick] (1.1,1.0) circle (4.2pt);
  \draw[<->,black] (0.46,0.59) -- (0.88,1.0);
  \node[font=\footnotesize,black!70] at (2.3,4.75) {松弛后的凸可行域};
  \node[font=\footnotesize,black!60] at (2.75,2.75) {凸包};
  \node[font=\footnotesize,black!70] at (5.15,0.5) {原问题的可行点};
  \draw[->,black!55] (4.0,0.68) -- (3.72,1.0);
  \node[font=\footnotesize,black!70] at (1.85,-0.5) {松弛最优值};
  \draw[->,black!55] (1.85,-0.33) -- (1.12,-0.19);
  \node[font=\footnotesize,black!70] at (3.45,-0.5) {原最优值};
  \draw[->,black!55] (3.45,-0.33) -- (2.38,-0.19);
\end{tikzpicture}

\emph{离散的可行点被一个凸集包住。带环的实心点是原问题的最优解，左下的空心圈是松弛的最优解，两条等值线之间的距离就是松弛间隙。}
\end{center}

图里那些实心点是原问题允许的方案，有限、离散、彼此孤立；虚线多边形是它们的凸包；外面那个椭圆是松弛之后实际能解的可行域。同一个线性目标沿等值线往左下推，在椭圆上比在点集上能推得更远，所以松弛的最优值一定不高于原问题的最优值。间隙有多大，图上是那两条虚线之间的距离；工程上就是你敢对交付方案打的包票。

\subsection{把``半正定''当成一条约束}

线性规划要求一个仿射向量落在非负象限里。半正定规划做的是同一件事，只是把``每个分量非负''换成``一个对称矩阵半正定''：

\[
\begin{aligned}
\text{minimize}\quad&c^{\trans}x\\
\text{subject to}\quad&F_0+\sum_{i=1}^{n}x_iF_i\succeq0,\\
&Ax=b,
\end{aligned}
\]

其中 \(F_0,\ldots,F_n\) 是给定的对称矩阵。左边那个矩阵对 \(x\) 是仿射的，所以约束是半正定锥在一个仿射映射下的原像，可行域凸——半正定锥本身为什么是凸锥，见\hyperref[sec:09-Convex-Optimization/02-Convex-Sets/03]{``从方向的集合到半正定锥''}。这种形式的约束通常叫\textbf{线性矩阵不等式}，简称 LMI。

线性规划是它的退化情形：把各个标量不等式摆在对角线上，要求一个对角矩阵半正定，等价于每个对角元非负。二阶锥约束也能嵌进来。所以 SDP 是一个比前几篇都宽的框子，宽到很多外观毫不相干的问题都落得进去。

宽在哪里？关键是矩阵不等式能一次性管住\emph{所有方向}。对称矩阵 \(A(x)\) 的最大特征值不超过 \(t\)，等价于

\[
tI-A(x)\succeq0,
\]

因为后者恰好是说 \(t-\lambda_i(A(x))\ge0\) 对每个特征值成立。同样，\(mI\preceq X\preceq MI\) 把 \(X\) 的整个谱压进区间 \([m,M]\)。这类条件没法用有限条逐元素不等式替代——谱是所有方向上二次型的极值，不是几个位置上的数。

Schur 补则负责把带逆矩阵的表达式拉平。若 \(P\succ0\)，

\[
\begin{bmatrix}P&x\\x^{\trans}&t\end{bmatrix}\succeq0
\quad\Longleftrightarrow\quad
t-x^{\trans}P^{-1}x\ge0 .
\]

左边对 \((P,x,t)\) 全体仿射，右边却含着 \(P^{-1}\)。于是二次型的上界约束变成了 LMI，求解时也不必显式形成逆矩阵。碰到分块矩阵的半正定约束，先问一句它是不是某个二次型或某个逆的上境图，往往就能认出隐藏的 SDP 结构。

\subsection{扔掉秩约束，就得到一个能解的下界}

回到机房划分。令 \(x\in\set{-1,+1}^{n}\)，代价可以写成 \(\tfrac12\trace(WX)\) 加常数，其中 \(W\) 是调用量矩阵，\(X=xx^{\trans}\)。矩阵 \(X\) 的三条性质是等价刻画：对角线全为 \(1\)，\(X\succeq0\)，以及 \(\rank(X)=1\)。前两条对 \(X\) 是凸的；第三条不是——秩不超过 \(r\) 的矩阵集合天生非凸，两个秩一矩阵相加一般是秩二。

上一篇的换坐标在这里无能为力。这个非凸性是问题的结构，不是坐标的产物。于是换个办法：把秩约束整个删掉，剩下

\[
\begin{aligned}
\text{minimize}\quad&\tfrac12\trace(WX)\\
\text{subject to}\quad&\diag(X)=\mathbf{1},\\
&X\succeq0 .
\end{aligned}
\]

这是一个标准 SDP，多项式时间可解。

\begin{proposition}
设原问题在集合 \(S\) 上最小化 \(f\)，松弛问题在 \(S'\supseteq S\) 上最小化同一个 \(f\)，则松弛的最优值不超过原问题的最优值。
\end{proposition}

骨架一步：原问题的每个可行点都是松弛问题的可行点，且目标值不变，所以松弛问题的可行值集合包含原问题的，取下确界只会更小或相等。条件对账：论证只用到``包含''与``目标函数相同''这两件事，没有用到 \(S'\) 凸，没有用到 \(f\) 凸，也没有用到 SDP 的任何性质。反过来说，这个论证同样没有对间隙的大小许下任何承诺——把 \(S'\) 取成整个空间，界照样成立，只是毫无用处。

所以有价值的松弛必须在两个方向上同时交代：它给的界要真的接近，而这件事需要单独证明。Max-Cut 的经典结果就是这样得来的——从 SDP 的解做 Cholesky 分解得到一组单位向量，用一个随机超平面把它们切成两边，再算这个随机方案的期望代价，才推出 \(0.878\) 这个比值。舍入与分析是额外的一整套论证，不是松弛自带的赠品，细节见\hyperref[sec:09-Convex-Optimization/09-Complexity-and-Approximation/02]{``凸松弛把 NP-hard 问题变成可解问题''}，问题难度本身的分类见\hyperref[sec:09-Convex-Optimization/09-Complexity-and-Approximation/01]{``P、NP 与 NP-hard：给问题按计算难度分类''}。

交付时因此要分三样东西报告：松弛给出的界、舍入后方案的实际代价、以及这个方案在原始约束下确实可行。只说``我解了一个 SDP''，等于什么都没说。

\subsection{同一个框子还装着协方差与控制}

SDP 的第二个入口是合法性。估计协方差矩阵时，逐元素地去拟合样本二阶矩，得到的东西未必是合法的协方差——只要某个方向上的二次型为负，就意味着某个线性组合的方差是负数。把 \(\Sigma\succeq0\) 写成约束，这件事从此不会发生。在此之上再加稀疏惩罚、低秩结构或谱区间，就得到协方差补全、实验设计和核学习里的一批凸模型；投资组合问题里的风险项也依赖同一个半正定性，见\hyperref[sec:09-Convex-Optimization/07-Applications/04]{``风险、收益与投资组合''}。

低秩这条线要小心。\(\set{X:\rank(X)\le r}\) 非凸，直接优化它会退回非凸的老路。常见替代是核范数 \(\norm{X}_*=\sum_i\sigma_i(X)\)，它是秩的凸代理，作用机制与 \(\ell_1\) 之于稀疏是同一套；或者干脆删掉秩约束做松弛，事后检查解出来的矩阵是不是接近低秩。传感器定位是个干净的例子：未知位置的 Gram 矩阵 \(G\) 满足 \(G\succeq0\)，两点距离的平方 \(G_{ii}+G_{jj}-2G_{ij}\) 对 \(G\) 是线性的，唯一麻烦的是``所有点必须落在三维空间里''这句话等于 \(\rank(G)\le3\)。删掉它得到 SDP；若解出的 \(G\) 秩明显高于 \(3\)，它可能把距离数据拟合得很好，却不对应任何真实布局，必须投影降维再回到原始距离上核验误差。

第三个入口在控制。线性系统 \(\dot z=Az\) 稳定，等价于存在 \(P\succ0\) 使 \(A^{\trans}P+PA\prec0\)。这是关于 \(P\) 的一组线性矩阵不等式，找 Lyapunov 函数因此变成解一个 SDP，而不是靠经验凑一个二次型，背景见\hyperref[sec:15-Differential-Equations-and-Dynamical-Systems/02-Stability-and-Dynamical-Systems/02]{``稳定性与 Lyapunov 函数''}。控制文献里大量``某某性能指标的 LMI 刻画''，走的都是这条路。

\subsection{表达力要付账}

框子宽不是白给的。\(n\times n\) 的对称矩阵变量有 \(n(n+1)/2\) 个自由度，通用内点法每一步要处理这么大的线性系统，矩阵一上千维，时间和内存都会顶到天花板。所以建模时的顺序应该反过来：先看问题是不是 LP，不是再看是不是 SOCP，实在需要谱约束或者需要一个松弛下界，才升级到 SDP。能用窄框子解决的事，不要用宽框子。

大规模问题另有出路——利用稀疏块结构、改用低秩参数化、或者换成一阶算子分裂方法。这些做法能把规模推上去，但会改变你最后能拿到的精度和证书质量。这是要事先想清楚的交换，不是实现细节。

到这里，这一单元把四类可解形式过了一遍：线性规划、二次规划、变换后现形的凸问题、半正定规划。它们共享同一句承诺——写成标准形式，就能拿回一个最优解和一份下界。可这份下界究竟是从哪儿冒出来的，前面一直在用，还没有造。下一单元从 Lagrangian 开始，把它造出来。

